\scanpage{88}
% loc 1–5
J'ai fait le calcul de $\rho$ (en termes de $b,c,b',c',t$) dans le cas où $b=\operatorname{Tr}u$ et $b'=\operatorname{Tr}v$ sont nuls, \deleted{on trouve}\inserted{et} supposant que $u,v$ plongés dans A Alg d'Az. On trouve alors
% loc 6–7
\[
\left.\begin{aligned}
\rho&=2\frac{t^2+4cc'}{t^2-4cc'}\\
\text{d'où}\quad \frac{\rho}{2}-1&=8\frac{cc'}{t^2-4cc'}
\end{aligned}\right|\ \text{si } b=b'=0
\]
% loc 6–7, right-hand NB
\margin{\uline{NB} $t^{\unsure{1|2}}-4cc'=\delta(uv)$ discriminant de $uv$}
% loc 8–10
On va en déduire le calcul dans le cas général. [Il suffit de le faire sur $\mathbb Z[\frac12][B,C,B',C',T]$.) . On trouve
% loc 11–12, display (11) and annotations beneath equalities
\[
\rho\underset{\substack{\downarrow\\\text{OK si 2 inv.}\\\text{dans }k}}{=}
2\frac{(2t-bb')^2+\delta(u)\delta(v)}{(2t-b\deleted{\unsure{b}}b')^2-\delta(u)\delta(v)}
\underset{\substack{|\\\text{OK}\\\text{ds ts les}\\\text{cas}}}{=}
\frac{2t^2-2bb't+\bigl(b^2b'^2-2(b^2c'+b'^2c)+8cc'\bigr)}
{t^2-bb't+(b^2c'+b'^2c-4cc')} \tag*{(11)}
\]
% loc 13, bracketed specialization
[ En \uline{car. 2}, cela donne $\displaystyle\rho\underset{\text{car. 2}}{=}\frac{(bb')^2}{t^2-bb't+(b^2c'+b'^2c)}$. ]
% loc 14–15
\[
\rho_0\overset{\unsure{\uparrow}}{=}\rho-2
=\frac{b^2b'^2-4(b^2c'+b'^2c)+16cc'}{t^2-bb't+(b^2c'+b'^2c-4cc')}
=\frac{A(b,b',c,c')}{B(t,b,b',c,c')\deleted{\unsure{t}}}
\tag*{\deleted{(\unsure{12|**})}}
\]
\margin{invariant anh. rectifié}
% loc 16–17
Donc
\[ \rho B(t,b,b',c,c')\deleted{\unsure{\rho}}-A(b,b',c,c')=0 \]
% archive sheet number: 87
